\chapter{積分の応用}\label{章:積分の応用}

定積分は,細かく分けた量を足し合わせる極限として定義した.この考えを使うと,曲線に囲まれた面積だけでなく,立体の体積や曲線の長さも求められる.ただし,どの問題にも被積分関数が最初から与えられるわけではない.細い部分を一つ取り出し,その面積,体積,長さをどの式で近似するかを決める必要がある.

その際に確かめたいことが三つある.一つは,分割した部分が対象全体を覆っているか.もう一つは,同じ部分を重複して数えていないか.そして,面積や道のりのような非負の量に,符号付きの差をそのまま使っていないかである.この章では,具体的な図形と運動を通して,それぞれの確認がどの式に現れるかを見ていく.

\section{面積}\label{節:曲線間の面積}

\sectionref{節:積分符号面積}では,曲線と $x$ 軸の間の面積を定積分で表した.では,2つの曲線に挟まれた部分の面積はどう求めればよいだろうか.

$a<b$ とし,\,$f,g$ は $[a,b]$ で連続で,\,$f(x)\ge g(x)$ を満たすとする.求めるのは
\[
a\le x\le b,\qquad g(x)\le y\le f(x)
\]
で表される部分の面積 $S$ である.この部分を $x$ 軸に垂直な細い帯に分け,各帯を長方形で近似する.長方形の面積の和の極限として,上下の差 $f(x)-g(x)$ の定積分に帰着させることを目標とする.

区間 $[a,b]$ を $n$ 等分し,\,$\Delta x=\frac{b-a}{n}$,\,$x_k=a+k\Delta x$ と置く.\,$k$ 番目の帯は,幅 $\Delta x$,高さ $f(x_k)-g(x_k)$ の長方形で近似できる.よって
\[
S\approx\sum_{k=1}^n\{f(x_k)-g(x_k)\}\Delta x
\]
である.右辺は連続関数 $f-g$ のリーマン和なので,\,$n\to\infty$ として
\begin{equation}\label{式:上下の差の面積}
S=\int_a^b\{f(x)-g(x)\}\dd x
\end{equation}
を得る.

長方形の高さは上の値から下の値を引いた差であり,曲線が $x$ 軸の上にあるかどうかは使っていない.実際,\,$f,g$ に同じ定数 $c$ を足しても,\,$(f+c)-(g+c)=f-g$ は変わらない.したがって,上下の曲線がともに $x$ 軸の下にあっても,同じ式が使えるのだ.

ここで1つ断りをいれておこう.式\eqref{式:上下の差の面積}は,区間全体で $f(x)\ge g(x)$ であることを前提にしている.2曲線が途中で交差すると,上下が入れ替わった区間では $f-g$ が負になり,面積を引いてしまう.1つの式で書けば
\[
S=\int_a^b|f(x)-g(x)|\dd x
\]
であるが,実際に計算するには,上下関係が変わる点で区間を分けて絶対値を外す.交点の $x$ 座標は区切りの候補であり,交点を求めただけで上下が決まるわけではない.差を因数分解するか,各区間で1つの値を代入して符号を確かめる.

\begin{bookpoint}
\textbf{POINT}\par
交点の $x$ 座標で区間を分け,各区間で(上の関数)$-$(下の関数)を積分して足す.
\end{bookpoint}

\begin{example}{交差する曲線の面積}
2曲線 $y=x^3-3x$ と $y=x$ に囲まれた部分の面積を求めよ.
\end{example}

\paragraphbf{方針}
交点で区間を分け,差 $x^3-4x$ の符号を調べて絶対値を外す.

\paragraphbf{解答}
$8$

\paragraphbf{解法}
交点の $x$ 座標は
\[
x^3-3x=x
\]
より,
\[
x(x+2)(x-2)=0
\]
であるから,\,$x=-2,0,2$ である.差 $x^3-4x=x(x+2)(x-2)$ は $-2<x<0$ で正,\,$0<x<2$ で負である.よって面積は
\[
S=\int_{-2}^0(x^3-4x)\dd x+\int_0^2(4x-x^3)\dd x
\]
となる.図\ref{図:面積の上下交換}では,左右の部分が原点について対称である.積分でも,第1項で $x=-t$ と置けば
\[
\int_{-2}^0(x^3-4x)\dd x=\int_0^2(4t-t^3)\dd t
\]
となり,第2項と一致する.したがって
\begin{align*}
S&=2\int_0^2(4x-x^3)\dd x\\
 &=2\left[2x^2-\frac{x^4}{4}\right]_0^2\\
 &=2(8-4)\\
 &=8.
\end{align*}

\begin{figure}[htbp]
\centering
\begin{tikzpicture}[x=1.6cm,y=0.65cm,>=Stealth]
\fill[blue!12]plot[domain=-2:0,samples=60](\x,{\x*\x*\x-3*\x})--plot[domain=0:-2,samples=30](\x,\x)--cycle;
\fill[orange!17]plot[domain=0:2,samples=30](\x,\x)--plot[domain=2:0,samples=60](\x,{\x*\x*\x-3*\x})--cycle;
\draw[->](-2.3,0)--(2.4,0)node[right]{$x$};
\draw[->](0,-2.8)--(0,3.4)node[above]{$y$};
\draw[blue!70!black,thick]plot[domain=-2.13:2.13,samples=120](\x,{\x*\x*\x-3*\x});
\draw[orange!80!black,thick](-2.2,-2.2)--(2.2,2.2);
\foreach \x in {-2,2}{\draw[dashed](\x,0)--(\x,\x);\node[below]at(\x,0){$\x$};}
\node[below right]at(0,0){$O$};
\node[blue!70!black]at(-1.25,2.8){$y=x^3-3x$};
\node[orange!80!black]at(1.6,3.0){$y=x$};
\end{tikzpicture}
\caption{同じ2曲線でも,交点の左右で上側の曲線が入れ替わる.}\label{図:面積の上下交換}
\end{figure}

区間を分けずに $\int_{-2}^2(x^3-4x)\dd x$ とすると,奇関数の積分なので $0$ になる.この $0$ は,左の部分の面積を正,右の部分の面積を負として足した値である.面積を求めるには,符号が変わる $x=0$ で区間を分けなければならない.

\subsection{積分変数の選択}

細い帯は縦に取る必要はない.\,$c\leq y\leq d$ において,右側の境界が $x=q(y)$,左側が $x=p(y)$ と表せるなら,帯の横の長さは $q(y)-p(y)$ なので
\begin{equation}\label{式:左右の差の面積}
S=\int_c^d\{q(y)-p(y)\}\dd y
\end{equation}
である.式\eqref{式:上下の差の面積}では上下の差,この式では左右の差を使う.変数の文字を交換するだけでなく,実際にどの向きに帯を取るかが変わっている.

たとえば $x=y^2$ と $x=2-y$ に囲まれる部分では,交点の $y$ 座標は $-2,1$ である.その間では直線が右側なので,
\[
S=\int_{-2}^1(2-y-y^2)\dd y
 =\left[2y-\frac{y^2}{2}-\frac{y^3}{3}\right]_{-2}^1
 =\frac92.
\]
縦に切ると放物線の上側と下側を $y=\pm\sqrt x$ に分け,さらに直線との関係も場合分けすることになる.横に切れば,一つの区間で一つの差になる.境界をどちらの変数で表すと少ない場合分けで済むかが,積分変数を選ぶ基準になる.



%
%
\section{断面積と体積}\label{節:断面積と体積}

立体を $x$ 軸に垂直な平面で切り,位置 $x$ における断面積を $A(x)$ とする.幅 $\Delta x$ の薄い部分を,底面積 $A(x)$,高さ $\Delta x$ の柱体で近似すれば,その体積は $A(x)\Delta x$ である.断面積が連続に変化する通常の立体では,分割を細かくした極限から
\begin{equation}\label{式:断面積の体積}
V=\int_a^b A(x)\dd x
\end{equation}
を得る.断面が途中で形を変える場合も,変わる位置で有限個の区間に分け,各区間で同じ考えを使える.

この式の $A(x)$ は長さではなく面積である.長さの単位を $\mathrm{cm}$ とすれば,\,$A(x)$ は $\mathrm{cm}^2$,幅 $\dd x$ を掛けた体積は $\mathrm{cm}^3$ になる.図形の式が合っているか迷ったとき,単位も確認の助けになる.

\begin{example}{正方形の断面}
平面上で $0\leq x\leq1$,\,$x^2\leq y\leq x$ となる部分を底面とする.各 $x$ において,底面内の線分を一辺とする正方形が,底面に垂直な断面になるような立体の体積を求めよ.
\end{example}

底面内の線分の長さは $x-x^2$ である.この長さをそのまま積分すると底面積になるが,今求めたいのは体積なので,正方形の断面積
\[
A(x)=(x-x^2)^2=x^2-2x^3+x^4
\]
を積分する.
\[
V=\int_0^1(x^2-2x^3+x^4)\dd x
 =\left[\frac{x^3}{3}-\frac{x^4}{2}+\frac{x^5}{5}\right]_0^1
 =\frac1{30}.
\]
同じ底面でも,断面が正三角形なら面積は $\frac{\sqrt3}{4}(x-x^2)^2$ であり,体積も上の値の $\frac{\sqrt3}{4}$ 倍になる.底面の情報だけでは体積は決まらず,断面の形が必要である.


%
\section{回転体}\label{節:回転体の体積}

$0\leq g(x)\leq f(x)$ の領域 $g(x)\leq y\leq f(x)$ を $x$ 軸のまわりに回すと,位置 $x$ の断面は外半径 $f(x)$,内半径 $g(x)$ の円環になる.よって
\begin{equation}\label{式:円環断面の体積}
V=\pi\int_a^b\{f(x)^2-g(x)^2\}\dd x.
\end{equation}
下側が軸そのもの,すなわち $g(x)=0$ なら断面は円板で,\,$V=\pi\int_a^b f(x)^2\dd x$ となる.

回転軸が $y=c$ なら,半径は $y$ 座標そのものではなく $|y-c|$ である.たとえば $c\leq g(x)\leq f(x)$ なら,外半径は $f(x)-c$,内半径は $g(x)-c$ になる.軸が領域より上にあるときは,下側の境界のほうが軸から遠い.上下の順序をそのまま外半径・内半径の順序にしてはいけない.

\subsection{回転後の重なり}

ここまでは,縦の線分が回転軸の片側にある場合を考えた.では,線分が軸をまたぐとどうなるだろうか.

縦の線分 $g(x)\le y\le f(x)$ が回転軸 $y=c$ をまたぐと,回転後の断面には中心まで点が入る.したがって内半径は $0$ であり,外半径は両端のうち軸から遠いほうの距離である.まとめると,
\[
R(x)=\max\{|f(x)-c|,|g(x)-c|\},
\]
\[
r(x)=
\begin{cases}
0,&g(x)\le c\le f(x),\\
\min\{|f(x)-c|,|g(x)-c|\},&\text{それ以外}
\end{cases}
\]
として,断面積は $\pi\{R(x)^2-r(x)^2\}$ になる.

軸より上の部分と下の部分を別々に回すと,同じ円板の部分を2度作ることがある.2つの体積を加えると,その部分を重複して数えてしまう.先に回転後の断面を考え,できた円板や円環を合わせた図形の面積を求める必要がある.

\begin{example}{軸をまたぐ領域を回す}
$0\le x\le2$,\,$x-1\le y\le1$ で表される領域を $x$ 軸のまわりに1回転してできる立体の体積を求めよ.
\end{example}

\paragraphbf{方針}
線分が軸をまたぐかどうかが変わる $x=1$ で区間を分け,円板と円環の断面積をそれぞれ積分する.

\paragraphbf{解答}
$\frac{5\pi}{3}$

\paragraphbf{解法}
$0\le x\le1$ では $x-1\le0\le1$ なので,線分は軸をまたぐ.上端の距離は $1$,下端の距離は $1-x\le1$ だから,断面は半径 $1$ の円板になる.\,$1\le x\le2$ では線分全体が軸の上にあり,断面は外半径 $1$,内半径 $x-1$ の円環になる.図\ref{図:回転断面の区別}は,この違いを断面で示している.よって
\begin{align*}
V&=\pi\int_0^1 1\dd x+\pi\int_1^2\{1-(x-1)^2\}\dd x\\
 &=\pi+\pi\left[x-\frac{(x-1)^3}{3}\right]_1^2\\
 &=\pi+\pi\left(2-\frac13-1\right)\\
 &=\frac{5\pi}{3}.
\end{align*}

\begin{figure}[htbp]
\centering
\begin{tikzpicture}[scale=1.0,>=Stealth]
\begin{scope}[xshift=-3.6cm]
\fill[blue!12](0,-1)--(0,1)--(2,1)--cycle;
\draw[->](-0.3,0)--(2.5,0)node[right]{$x$};
\draw[->](0,-1.3)--(0,1.5)node[above]{$y$};
\draw[thick](0,1)--(2,1)--(0,-1);
\draw[blue!70!black,very thick](0.5,-0.5)--(0.5,1);
\draw[orange!85!black,very thick](1.5,0.5)--(1.5,1);
\node[below]at(1,0){$1$};\node[below]at(2,0){$2$};
\node[below]at(1,-1.7){回す前の領域};
\end{scope}
\begin{scope}[xshift=0.1cm]
\fill[blue!12](0,0)circle(1);
\draw[blue!70!black,thick](0,0)circle(1);
\draw[->](0,0)--(0.707,0.707)node[midway,left]{$1$};
\fill(0,0)circle(0.025);
\node[below]at(0,-1.45){$x=\frac12$：円板};
\end{scope}
\begin{scope}[xshift=3.2cm]
\fill[orange!17,even odd rule](0,0)circle(1)(0,0)circle(0.5);
\draw[orange!85!black,thick](0,0)circle(1);
\draw[orange!85!black,thick](0,0)circle(0.5);
\draw[->](0,0)--(0.707,0.707)node[midway,right]{$1$};
\draw[->](0,0)--(-0.5,0)node[midway,above]{$\frac12$};
\fill(0,0)circle(0.025);
\node[below]at(0,-1.45){$x=\frac32$：円環};
\end{scope}
\end{tikzpicture}
\caption{軸を含む線分は円板を作り,軸から離れた線分は円環を作る.}\label{図:回転断面の区別}
\end{figure}

最初から $\pi\int_0^2\{1-(x-1)^2\}\dd x$ とすると,\,$x<1$ でも中心の半径 $1-x$ の円板をくり抜いてしまう.しかし,元の線分は軸を含んでいるので,回転後にその穴はない.

1つの $x$ に対して元の領域が複数の線分に分かれる場合には,各線分が回転して覆う半径の範囲を調べる.範囲が重なるなら,その部分を1度だけ数える.元の平面図形では重なっていない2つの部分を回す場合にも,この確認が必要になる.
%
\subsection{円筒による分割}

同じ立体でも,回転軸に垂直に切る代わりに,軸に平行な薄い帯を回す方法がある.たとえば $0\leq a\leq x\leq b$,\,$g(x)\leq y\leq f(x)$ を $y$ 軸のまわりに回す.幅 $\Delta x$ の帯は,半径がほぼ $x$,高さがほぼ $f(x)-g(x)$ の薄い円筒の層になる.

半径 $x$ と $x+\Delta x$ の間の円環の面積は
\[
\pi\{(x+\Delta x)^2-x^2\}
 =2\pi x\Delta x+\pi(\Delta x)^2
\]
である.分割を細かくすると二次の項による誤差の和は零に近づくので,体積は
\begin{equation}\label{式:円筒殻の体積}
V=2\pi\int_a^b x\{f(x)-g(x)\}\dd x
\end{equation}
となる.これを円筒殻による方法という.\,$2\pi x$ は円周の長さ,\,$f-g$ は層の高さである.

\begin{example}{体積の別解}
$0\leq x\leq1$,\,$x^2\leq y\leq x$ の領域を $y$ 軸のまわりに回転させた立体の体積を求めよ.
\end{example}

縦の帯を回すと,半径 $x$,高さ $x-x^2$ の円筒殻になるから,
\[
V=2\pi\int_0^1x(x-x^2)\dd x
 =2\pi\left[\frac{x^3}{3}-\frac{x^4}{4}\right]_0^1
 =\frac{\pi}{6}.
\]
横に切る方法でも確かめられる.\,$0\leq y\leq1$ で,領域の左右の端は $x=y$ と $x=\sqrt y$ である.したがって円環の外半径は $\sqrt y$,内半径は $y$ となり,
\[
V=\pi\int_0^1\{(\sqrt y)^2-y^2\}\dd y
 =\pi\left[\frac{y^2}{2}-\frac{y^3}{3}\right]_0^1
 =\frac{\pi}{6}.
\]
断面を円環として積み重ねるか,円筒の層として重ねるかによって,被積分関数は異なる.しかし,それぞれが立体を重複なく分けているので同じ体積になる.

式\eqref{式:円筒殻の体積}で $x\geq0$ としたのは,異なる $x$ が異なる半径に対応するためである.領域が $y$ 軸をまたぐと,\,$x$ と $-x$ が同じ半径になる.たとえば $-1\leq x\leq1$,\,$0\leq y\leq1$ の長方形を回すと半径 $1$,高さ $1$ の円柱で,体積は $\pi$ である.ところが $2\pi\int_{-1}^1|x|\dd x=2\pi$ は同じ層を二度数えている.絶対値で半径を正にしただけでは,重複は解消しない.



%
%
\section{曲線の長さ}\label{節:曲線の弧長}

面積は長方形で,体積は柱体で近似した.では,曲線 $y=f(x)$ の長さはどう求めればよいだろうか.長方形は面積を測るための形なので,長さの近似には使えない.

すぐに思いつく候補を2つ挙げて,外しておこう.1つは横の幅 $\int_a^b\dd x=b-a$ だが,これは横方向の移動しか数えていない.もう1つは縦方向の変化の総量 $\int_a^b|f'(x)|\dd x$ だが,水平な線分では $f'=0$ となり,長さ $b-a$ を数えられない.横と縦の両方の移動を同時に数える必要がある.そこで,曲線上の近い2点を線分で結び,折れ線の長さで近似しよう.

実験として,\,$y=\frac23x^{\frac32}$ の $0\le x\le3$ の部分を考える.区間を $n$ 等分し,分点に対応する曲線上の点を順に線分で結んだ折れ線の長さを $L_n$ とする.三平方の定理で各線分の長さを求めて足すと,次の表のようになる.
\begin{center}
\begin{tabular}{ccccccc}
\toprule
$n$ & $1$ & $2$ & $4$ & $8$ & $16$ & $32$\\
\midrule
$L_n$ & $4.58258$ & $4.63181$ & $4.65396$ & $4.66245$ & $4.66535$ & $4.66627$\\
\bottomrule
\end{tabular}
\end{center}
分点を増やすと折れ線は曲線に近づき,\,$L_n$ は増えながら $4.667$ 付近に近づいていく.一方,横の幅は $3$,縦の変化の総量は $2\sqrt3\approx3.464$ であり,どちらも折れ線より短い.この極限の値を積分で表そう.

$f$ は $[a,b]$ で連続微分可能,すなわち導関数 $f'$ も連続であるとする.\,$[a,b]$ を $n$ 等分し,\,$\Delta x=\frac{b-a}{n}$,\,$x_k=a+k\Delta x$,\,$\Delta y_k=f(x_k)-f(x_{k-1})$ と置く.\,$k$ 番目の線分の長さは
\[
\sqrt{(\Delta x)^2+(\Delta y_k)^2}
 =\sqrt{1+\left(\frac{\Delta y_k}{\Delta x}\right)^2}\,\Delta x
\]
である.平均値の定理により,\,$x_{k-1}<\xi_k<x_k$ を満たすある $\xi_k$ で $\frac{\Delta y_k}{\Delta x}=f'(\xi_k)$ となるから,折れ線の長さは
\[
\sum_{k=1}^n\sqrt{1+f'(\xi_k)^2}\,\Delta x
\]
と書ける.これは連続関数 $\sqrt{1+f'(x)^2}$ のリーマン和なので,\,$n\to\infty$ として
\begin{equation}\label{式:グラフの弧長}
L=\int_a^b\sqrt{1+f'(x)^2}\dd x
\end{equation}
を得る.これを曲線の\textbf{弧長}という.根号の中の $1$ が横方向の移動を,\,$f'(x)^2$ が縦方向の変化を数えている.有限個の点で折れ曲がっていても,各部分が連続微分可能なら,部分ごとの長さを加えればよい.

実験の曲線に当てはめてみよう.\,$f(x)=\frac23x^{\frac32}$ より $f'(x)=\sqrt x$ であるから,
\begin{align*}
L&=\int_0^3\sqrt{1+x}\dd x\\
 &=\left[\frac23(1+x)^{\frac32}\right]_0^3\\
 &=\frac23(8-1)\\
 &=\frac{14}{3}
\end{align*}
となる.\,$\frac{14}{3}=4.666\cdots$ であり,表の値が近づいていく先と一致するのだ.

\begin{bookpoint}
\textbf{POINT}\par
$f'(x)$ を求めて $1+f'(x)^2$ を整理し,\,$\int_a^b\sqrt{1+f'(x)^2}\dd x$ を計算する.
\end{bookpoint}

\subsection{媒介変数表示}

円や,縦向きの接線を持つ曲線は,1つの関数 $y=f(x)$ のグラフとして扱いにくい.そこで
\[
x=x(t),\qquad y=y(t),\qquad \alpha\leq t\leq\beta
\]
と表された曲線の長さを考えよう.\,$x(t),y(t)$ は連続微分可能であるとする.

$t$ を $\Delta t$ だけ増やしたとき,横の変化はほぼ $x'(t)\Delta t$,縦の変化はほぼ $y'(t)\Delta t$ である.よって,対応する線分の長さはほぼ
\[
\sqrt{\{x'(t)\Delta t\}^2+\{y'(t)\Delta t\}^2}
 =\sqrt{x'(t)^2+y'(t)^2}\,\Delta t
\]
となる.これを足し合わせた折れ線の極限として
\begin{equation}\label{式:媒介曲線の弧長}
L=\int_\alpha^\beta\sqrt{x'(t)^2+y'(t)^2}\dd t
\end{equation}
を得る.導関数が連続であることから近似の誤差は区間全体で小さくできるが,ここではその評価の詳細は省き,この式を用いる.途中で $\sqrt{x'(t)^2+y'(t)^2}$ が $0$ になる点があっても,この式はそのまま使える.

ここで1つ断りをいれておこう.この積分は,\,$t$ に沿って点が進んだ長さを数える.同じ部分を2度たどれば2度数える.たとえば $x=\cos t$,\,$y=\sin t$ を $0\leq t\leq4\pi$ で積分すると $4\pi$ になるが,図形としての単位円の周の長さは $2\pi$ である.図形そのものの長さを求めるなら,1度だけたどる区間を選ぶ.

媒介変数表示は面積にも使える.\,$x(t)$ が単調増加し,\,$y(t)\geq0$ であるとする.このとき曲線は $x$ の関数のグラフであり,曲線と $x$ 軸の間の面積は $\int_{x(\alpha)}^{x(\beta)}y\dd x$ である.\,$x=x(t)$ と置換すると $\dd x=x'(t)\dd t$ なので,
\begin{equation}\label{式:媒介曲線の面積}
S=\int_\alpha^\beta y(t)x'(t)\dd t
\end{equation}
となる.\,$x(t)$ が途中で折り返すなら,区間を分けて,曲線が囲む部分を確認する.この式に絶対値を付けるだけで,いつでも図形の面積になるわけではない.縦の帯が重なれば,同じ面積を2度数えるからである.

\begin{bookpoint}
\textbf{POINT}\par
長さは $\sqrt{x'(t)^2+y'(t)^2}$ を,面積は $y(t)x'(t)$ を $t$ で積分する.根号を外すときは $t$ の範囲で符号を確かめる.
\end{bookpoint}

\begin{example}{媒介曲線の面積と長さ}
$a>0$ とする.曲線
\[
x=a(t-\sin t),\qquad y=a(1-\cos t),\qquad 0\leq t\leq2\pi
\]
の長さと,この曲線と $x$ 軸に囲まれた部分の面積を求めよ.
\end{example}

\paragraphbf{方針}
長さは式\eqref{式:媒介曲線の弧長}で求め,根号の中を半角の公式で2乗の形にする.面積は $x(t)$ が単調増加であることを確かめてから式\eqref{式:媒介曲線の面積}を使う.

\paragraphbf{解答}
長さ $8a$,面積 $3\pi a^2$

\paragraphbf{解法}
$x'(t)=a(1-\cos t)$,\,$y'(t)=a\sin t$ である.

\textbf{長さ}\quad 根号の中は
\begin{align*}
x'(t)^2+y'(t)^2
 &=a^2\{(1-\cos t)^2+\sin^2t\}\\
 &=a^2(2-2\cos t)\\
 &=4a^2\sin^2\frac t2
\end{align*}
となる.平方根は $2a\left|\sin\frac t2\right|$ である.\,$0\leq t\leq2\pi$ では $0\leq\frac t2\leq\pi$ なので $\sin\frac t2\geq0$ であり,
\begin{align*}
L&=2a\int_0^{2\pi}\sin\frac t2\dd t\\
 &=\left[-4a\cos\frac t2\right]_0^{2\pi}\\
 &=4a-(-4a)\\
 &=8a.
\end{align*}

\textbf{面積}\quad $x'(t)=a(1-\cos t)$ は $0<t<2\pi$ で正なので,\,$x$ は区間全体で単調増加する.また $y\geq0$ で,両端では $y=0$ である.よって
\begin{align*}
S&=\int_0^{2\pi}a(1-\cos t)\cdot a(1-\cos t)\dd t\\
 &=a^2\int_0^{2\pi}(1-2\cos t+\cos^2t)\dd t\\
 &=a^2\int_0^{2\pi}\left(\frac32-2\cos t+\frac12\cos2t\right)\dd t\\
 &=a^2\cdot\frac32\cdot2\pi\\
 &=3\pi a^2.
\end{align*}
$\cos t$ と $\cos2t$ の項は,\,$0$ から $2\pi$ まで積分すると $0$ になる.
%
%
\section{位置と道のり}\label{節:運動と積分}

\sectionref{節:符号と道のり}では,速度が一定の区間をつないだ運動で,速度に時間を掛けて足すと位置の変化,速さに時間を掛けて足すと道のりになることを見た.では,速度が連続的に変わる場合はどうなるだろうか.

一直線上の運動で,時刻 $t$ の位置を $s(t)$,速度を $v(t)=s'(t)$,加速度を $a(t)=v'(t)$ とし,\,$v$ は連続であるとする.時刻 $t_0$ から $t_1$ までを細かく分けると,各小区間では速度がほぼ一定とみなせる.長さ $\Delta t$ の小区間で,位置の変化はほぼ $v(t)\Delta t$,進んだ距離はほぼ $|v(t)|\Delta t$ である.これらの和の極限として,位置の変化も道のりも定積分に帰着する.

位置の変化については,微積分の基本定理から直接
\[
s(t_1)-s(t_0)=\int_{t_0}^{t_1}v(t)\dd t
\]
を得る.左辺は位置の差,すなわち\textbf{変位}である.出発点に戻れば変位は $0$ になるが,途中で移動した距離が $0$ という意味ではない.

速度が正なら正の向きに,負なら負の向きに進む.向きによらず移動した距離を足すには,速さ $|v(t)|$ を積分して
\begin{equation}\label{式:道のり}
\text{道のり}=\int_{t_0}^{t_1}|v(t)|\dd t
\end{equation}
とする.速度には向きに対応する符号があり,速さは負にならない.この違いは,符号付きの定積分と面積の違いと同じである.

速度が $0$ になる時刻は,進む向きが変わる候補である.ただし,\,$0$ になるだけでは向きが変わるとは限らない.たとえば $v(t)=(t-1)^2$ は $t=1$ で $0$ になるが,前後で負にならない.道のりを計算するときは,速度が $0$ になる時刻で区間を分けたうえで,各区間の符号を確かめる.

\begin{bookpoint}
\textbf{POINT}\par
速度が $0$ になる時刻で区間を分け,各区間で $\int v\dd t$(位置の差)の絶対値を求めて足す.
\end{bookpoint}

\begin{example}{変位と道のり}
時刻 $0\le t\le4$ において,速度が $v(t)=(t-1)(t-3)$ で与えられる.初期位置は $s(0)=0$ とする.\,$t=4$ の位置,道のり,平均速度,平均の速さを求めよ.
\end{example}

\paragraphbf{方針}
速度を積分して位置 $s(t)$ を求める.\,$v(t)=(t-1)(t-3)$ は $1<t<3$ で負,それ以外では $0$ 以上なので,道のりは $[0,1]$,\,$[1,3]$,\,$[3,4]$ に分けて,各区間の位置の差の絶対値を足す.

\paragraphbf{解答}
位置 $\frac43$,道のり $4$,平均速度 $\frac13$,平均の速さ $1$

\paragraphbf{解法}
$v(t)=t^2-4t+3$ を積分し,\,$s(0)=0$ を使うと
\[
s(t)=\frac{t^3}{3}-2t^2+3t
\]
である.よって
\begin{align*}
s(4)&=\frac{64}{3}-32+12\\
 &=\frac43.
\end{align*}
次に速度の符号を調べる.\,$v$ は $0\le t<1$ で正,\,$1<t<3$ で負,\,$3<t\le4$ で正である.向きを変える時刻の位置は
\[
s(1)=\frac43,\qquad s(3)=0
\]
なので,道のりは
\begin{align*}
&\left|s(1)-s(0)\right|+\left|s(3)-s(1)\right|+\left|s(4)-s(3)\right|\\
&=\frac43+\frac43+\frac43\\
&=4
\end{align*}
となる.これは $\int_0^1v\dd t-\int_1^3v\dd t+\int_3^4v\dd t$ を計算したことと同じである.全時間は $4$ だから,平均速度と平均の速さはそれぞれ
\[
\frac{s(4)-s(0)}{4}=\frac13,\qquad
\frac{4}{4}=1
\]
となる.

\paragraphbf{速さが増える区間}
例題の加速度は $a(t)=2t-4$ である.加速度が正だからといって,速さが増えるとは限らない.\,$v(t)\ne0$ の時刻では
\[
\frac{\dd}{\dd t}|v(t)|=
\begin{cases}
a(t),&v(t)>0,\\
-a(t),&v(t)<0
\end{cases}
\]
であるから,速度と加速度が同符号なら速さが増え,異符号なら減る.例題では,\,$2<t<3$ で加速度は正だが速度は負なので,速さは減っている.逆に $1<t<2$ では速度も加速度も負なので,負の向きに進みながら速くなっているのだ.速さが増えるのは $1<t<2$ と $3<t<4$ である.

\paragraphbf{平面上の運動}
平面上を動く点 $\bigl(x(t),y(t)\bigr)$ では,速度は2つの成分 $\bigl(x'(t),y'(t)\bigr)$ を持ち,速さは三平方の定理から $\sqrt{x'(t)^2+y'(t)^2}$ となる.この速さを時刻 $\alpha$ から $\beta$ まで積分した道のりは,式\eqref{式:媒介曲線の弧長}と同じ式である.つまり,媒介変数 $t$ を時刻とみれば,曲線の長さの公式は動く点の道のりを求める式になっている.ただし運動では,同じ場所を再び通った分も道のりに含める.曲線を重複してたどった長さは,図形の長さとしては数えすぎだが,道のりとしては数えるべき量なのである.
%
\section{平均値}\label{節:関数の平均}

\sectionref{節:定積分性質}で導入した関数の平均値を,温度のように時間とともに連続的に変わる量に使ってみよう.その前に,平均を積分で表す理由を振り返っておく.

$n$ 個の数 $u_1,\ldots,u_n$ の平均は,和を個数で割ったものである.関数 $f$ の区間 $[a,b]$ における平均も,区間から等間隔に $n$ 個の値を取り,その平均の極限として考える.\,$\Delta x=\frac{b-a}{n}$ と置けば,\,$\frac1n=\frac{\Delta x}{b-a}$ であるから,
\[
\frac1n\sum_{k=1}^n f(a+k\Delta x)
 =\frac{1}{b-a}\sum_{k=1}^n f(a+k\Delta x)\Delta x
\]
となる.右辺の和はリーマン和なので,\,$f$ が連続なら,\,$n\to\infty$ として
\begin{equation}\label{式:関数の平均値}
\overline f=\frac{1}{b-a}\int_a^b f(x)\dd x
\end{equation}
を得る.\sectionref{節:定積分性質}ではこれを $f_{\mathrm{av}}$ と書いたが,本章では $\overline f$ と書く.\,$f\geq0$ の場合,底辺 $b-a$,高さ $\overline f$ の長方形は,曲線の下の部分と同じ面積になる.\,$f$ が負の値をとっても,平均は同じ式で定める.

\begin{example}{変化する温度の平均}
時刻 $0\leq t\leq12$ の温度を $T(t)=20+6\sin\frac{\pi t}{12}$ とする.この時間帯の平均温度を求めよ.
\end{example}

\paragraphbf{方針}
式\eqref{式:関数の平均値}に従って,\,$T(t)$ を $0$ から $12$ まで積分し,時間の長さ $12$ で割る.

\paragraphbf{解答}
$20+\frac{12}{\pi}$

\paragraphbf{解法}
\begin{align*}
\overline T
 &=\frac1{12}\int_0^{12}\left(20+6\sin\frac{\pi t}{12}\right)\dd t\\
 &=\frac1{12}\left[20t-\frac{72}{\pi}\cos\frac{\pi t}{12}\right]_0^{12}\\
 &=\frac1{12}\left\{240-\frac{72}{\pi}(-1-1)\right\}\\
 &=20+\frac{12}{\pi}.
\end{align*}
始めと終わりの温度はともに $20$ だが,途中の温度は $20$ より高い.両端の温度の平均 $20$ は,時間全体の平均にはならないのである.

\paragraphbf{平均値をとる点}
\sectionref{節:定積分性質}で見たように,連続な $f$ の最小値を $m$,最大値を $M$ とすると $m\leq\overline f\leq M$ であり,ある $c\in[a,b]$ で $f(c)=\overline f$ となる(積分の平均値の定理).ただし,その $c$ が区間の中央にあるとは限らず,1つとも限らない.上の例題では $\sin\frac{\pi t}{12}=\frac2\pi$ を満たす $t$ が $0<t<12$ に2つあり,どちらも中央の $t=6$ とは異なる.

\paragraphbf{重み付き平均}
各点を均等に数えない平均もある.非負の連続関数 $w$ が $\int_a^b w(x)\dd x>0$ を満たすとき,\,$w$ を重みとする\textbf{重み付き平均}を
\[
\frac{\int_a^b f(x)w(x)\dd x}{\int_a^b w(x)\dd x}
\]
で定める.分母は重みの総量であり,\,$w=1$ とすれば式\eqref{式:関数の平均値}に戻る.次節で扱う重心は,位置の座標を面積や長さを重みとして平均したものである.
%
\section{重心}\label{節:面積重心と線重心}

前節の重み付き平均を,密度が一様な薄い板の重心に使ってみよう.板を小さな部分に分け,\,$i$ 番目の部分の面積を $\Delta A_i$,位置をほぼ $(x_i,y_i)$ とする.密度が一様なら各部分の質量は面積に比例し,共通の密度は分子と分母で約分できる.よって重心の座標は,面積を重みとした平均
\[
\overline x\approx\frac{\sum_i x_i\Delta A_i}{\sum_i\Delta A_i},\qquad
\overline y\approx\frac{\sum_i y_i\Delta A_i}{\sum_i\Delta A_i}
\]
で表される.小さな部分として縦の細い帯をとり,この和を定積分に帰着させよう.

$a<b$ とし,\,$f,g$ は $[a,b]$ で連続で $g(x)\le f(x)$ を満たすとする.板が
\[
a\le x\le b,\qquad g(x)\le y\le f(x)
\]
で表され,その面積 $A$ は正であるとする.位置 $x$ の幅 $\Delta x$ の帯は,面積がほぼ $\{f(x)-g(x)\}\Delta x$ の長方形で,その中心は $\left(x,\frac{f(x)+g(x)}{2}\right)$ である.\,$y$ 座標の重み付きの和を作ると,帯1つ分は
\[
\frac{f(x)+g(x)}{2}\cdot\{f(x)-g(x)\}\Delta x
 =\frac12\{f(x)^2-g(x)^2\}\Delta x
\]
となる.分割を細かくすると
\begin{equation}\label{式:面積の重心}
\begin{aligned}
A&=\int_a^b(f-g)\dd x,\\
\overline x&=\frac1A\int_a^b x(f-g)\dd x,\\
\overline y&=\frac1{2A}\int_a^b(f^2-g^2)\dd x
\end{aligned}
\end{equation}
を得る.これを\textbf{面積重心}という.\,$\overline y$ を計算するとき,上端の高さ $f$ だけを帯の面積に掛けてはいけない.帯の重心は上端ではなく,上下の中央にある.

曲線に沿って曲げた一様な針金では,質量が比例するのは面積ではなく長さである.曲線 $C$ の長さを $L>0$ とし,曲線の小さな部分の長さを $\dd s$ で表すと,重心は
\begin{equation}\label{式:線の重心}
\overline x=\frac1L\int_C x\dd s,\qquad
\overline y=\frac1L\int_C y\dd s
\end{equation}
となる.これを\textbf{線重心}という.\,$\int_C y\dd s$ は新しい種類の積分ではなく,曲線を媒介変数で表せば,式\eqref{式:媒介曲線の弧長}と同じく
\[
\int_C y\dd s
 =\int_\alpha^\beta y(t)\sqrt{x'(t)^2+y'(t)^2}\dd t
\]
という普通の定積分で計算する.針金の重心を求めるときは,曲線を1度だけたどる区間を選ぶ.

\begin{bookpoint}
\textbf{POINT}\par
板では縦の帯の面積 $(f-g)\dd x$ を重みにして,帯の中心 $\left(x,\frac{f+g}{2}\right)$ の座標を平均する.針金では長さ $\dd s$ を重みにする.
\end{bookpoint}

\paragraphbf{半円板と半円弧}
半径 $a$ の上半円板と,その円弧だけでできた針金の重心を比べよう.どちらも $y$ 軸について対称なので $\overline x=0$ である.

半円板では $A=\frac{\pi a^2}{2}$,\,$g(x)=0$,\,$f(x)=\sqrt{a^2-x^2}$ である.式\eqref{式:面積の重心}の $\overline y$ の積分は
\begin{align*}
\frac12\int_{-a}^a(a^2-x^2)\dd x
 &=\frac12\left[a^2x-\frac{x^3}{3}\right]_{-a}^a\\
 &=\frac{2a^3}{3}
\end{align*}
となる.これを $A$ で割って,半円板の重心の高さは
\begin{equation}\label{式:半円板の重心}
\overline y_{\mathrm{面}}=\frac{2a^3}{3}\cdot\frac{2}{\pi a^2}=\frac{4a}{3\pi}
\end{equation}
である.

円弧は $x=a\cos t$,\,$y=a\sin t$,\,$0\le t\le\pi$ と表せる.\,$\sqrt{x'(t)^2+y'(t)^2}=a$ より $\dd s=a\dd t$,\,$L=\pi a$ であるから,
\begin{equation}\label{式:半円弧の重心}
\overline y_{\mathrm{線}}
 =\frac{1}{\pi a}\int_0^\pi(a\sin t)a\dd t
 =\frac{2a}{\pi}
\end{equation}
となる.\,$\frac{2a}{\pi}\approx0.64a$,\,$\frac{4a}{3\pi}\approx0.42a$ なので,円弧だけのほうが重心は高い.円板には直径の近くにも面積があるが,円弧だけの針金にはその部分がないからである.
%
\section{パップス＝ギュルダンの定理}\label{節:パップスの定理}

回転体の体積を円筒殻による方法で求めると,軸からの距離に帯の面積を掛けた積分が現れた.一方,重心も座標に帯の面積を掛けた積分で求めた.この2つは同じ積分ではないだろうか.そうなら,回転体の体積は面積と重心だけで書けるはずである.

回転軸を $x$ 軸とし,図形が $y\ge0$ の側にあるとする.図形を $x$ 軸に平行な横の細い帯に分け,高さ $y$ の帯の横の長さを $h(y)$ とする.帯が複数に分かれるなら,\,$h(y)$ はそれらの長さの和とする.図形の高さの範囲を $c\le y\le d$ としよう.

横の帯を $x$ 軸のまわりに回すと,半径 $y$,軸に沿う長さ $h(y)$ の円筒の層になる.式\eqref{式:円筒殻の体積}で $x$ と $y$ の役割を入れ替えれば
\[
V=2\pi\int_c^d y\,h(y)\dd y
\]
である.他方,横の帯の中心の高さは $y$ であるから,面積 $A$ と面積重心の高さ $\overline y$ は
\[
A=\int_c^d h(y)\dd y,\qquad
A\overline y=\int_c^d y\,h(y)\dd y
\]
を満たす.

\sectionref{節:面積重心と線重心}では,\,$A\overline y$ を縦の帯から $\frac12\int_a^b(f^2-g^2)\dd x$ として求めた.横の帯から出した上の積分と同じ値になるのだろうか.図形を縦横の線で細かい長方形に分け,各長方形の面積 $\Delta x\Delta y$ に高さ $y$ を掛けた和 $\sum y\,\Delta x\Delta y$ を考えよう.同じ縦の列にある長方形を先に足すと,列1つ分は $\int_{g(x)}^{f(x)}y\dd y\cdot\Delta x=\frac12\{f(x)^2-g(x)^2\}\Delta x$ に近く,これは縦の帯1本分の寄与である.同じ横の行を先に足すと,行1つ分は $y\,h(y)\Delta y$ に近く,これは横の帯1本分の寄与である.足す順序を変えても和は変わらないので,分割を細かくした極限で2つの積分は一致する.どちらも,図形全体で(高さ)$\times$(面積)を足した同じ量を,分け方を変えて計算しているのだ.厳密な証明には極限の扱いに注意が要るが,ここでは立ち入らない.実際,半径 $a$ の上半円板では $h(y)=2\sqrt{a^2-y^2}$ であり,
\[
\int_0^a 2y\sqrt{a^2-y^2}\dd y
 =\left[-\frac23(a^2-y^2)^{\frac32}\right]_0^a
 =\frac{2a^3}{3}
\]
となって,\sectionref{節:面積重心と線重心}で縦の帯から求めた値と一致する.

$V$ の式と $A\overline y$ の式に,同じ積分 $\int_c^d y\,h(y)\dd y$ が現れている.これを $A\overline y$ に置き換えれば,
\[
V=2\pi\overline yA
\]
を得るのだ.

\begin{theorem}{パップス＝ギュルダンの体積に関する定理}
一様な平面図形の面積を $A$,その面積重心と,図形と同じ平面内にある回転軸との距離を $R$ とする.図形が軸の片側にあり,回転によって内部を重複して数えない場合,できる回転体の体積は
\[
V=2\pi R A
\]
である.すなわち,面積に,重心が1回転するときに進む円周の長さを掛ければよい.
\end{theorem}

上の証明では,境界が区分的に滑らかで,断面積による通常の積分ができる図形を考えている.軸が境界に触れても内部を横切らなければ,上の帯による証明がそのまま使える.軸が内部を横切る場合には,軸の両側が同じ場所を回転して重なり得るので,この定理は使えない.重心が軸上にあるから体積が $0$ になる,という使い方はできないのである.

\subsection{回転面の面積}

体積では面積重心が現れた.では,曲線を回してできる面の面積ではどうなるだろうか.結論から言えば,今度は線重心が現れる.その理由を,回転面の小さな部分から考えよう.

回転軸と同じ平面内にあり,軸の片側にある曲線を考える.曲線を短い線分の折れ線で近似し,各線分が回ってできる面の面積を足し合わせて,定積分に帰着させる.

軸からの距離が $r_1,r_2$ の2点を結ぶ長さ $\Delta s$ の線分を軸のまわりに回すと,一般には円すい台の側面ができる.軸に平行な線分なら円筒面,軸に垂直な線分なら円環または円板になる.いずれの場合も,その面積は
\[
\pi(r_1+r_2)\Delta s
\]
である.円すい台の場合は,大きい円すいの側面積から切り取った円すいの側面積を引き,相似比を使って整理するとこの式になる.円筒面では $r_1=r_2$ として $2\pi r_1\Delta s$,円環では $\Delta s=|r_2-r_1|$ として $\pi|r_2^2-r_1^2|$ であり,どちらもこの式に一致する.線分を十分短くすれば $r_1,r_2$ はほぼ同じ値 $r$ になるので,面積はほぼ $2\pi r\Delta s$ である.

曲線が区分的に連続微分可能で,回転によって同じ面を重複して作らないとき,折れ線近似の極限として,回転面の面積は
\begin{equation}\label{式:回転面の面積}
S=2\pi\int_C r\dd s
\end{equation}
になる.ここでは,この極限の正当化の詳細は省く.たとえば $y=f(x)\geq0$($a\leq x\leq b$)を $x$ 軸のまわりに回すなら,\,$r=f(x)$,\,$\dd s=\sqrt{1+f'(x)^2}\dd x$ であるから
\[
S=2\pi\int_a^b f(x)\sqrt{1+f'(x)^2}\dd x
\]
となる.得られるのは,指定した曲線が作る面の面積である.立体の両端を塞ぐ円板などは,その曲線に含めていなければ別に数える.

\begin{theorem}{パップス＝ギュルダンの面積に関する定理}
一様な平面曲線の長さを $L$,その線重心と,曲線と同じ平面内にある回転軸との距離を $R$ とする.曲線が軸の片側にあり,1回転によって同じ面を重複して作らない場合,曲線が作る回転面の面積は
\[
S=2\pi R L
\]
である.
\end{theorem}

実際,軸を $x$ 軸,曲線を $y\geq0$ の側にとれば $r=y$ である.式\eqref{式:線の重心}より $\int_C y\dd s=L\overline y=LR$ となるので,式\eqref{式:回転面の面積}から $S=2\pi RL$ が従う.回転軸の位置や向きが違っていても,軸に合わせて座標をとり直せば同じ議論になる.

2つの定理で掛ける円周は,図形の端や図形の中心が描く円周ではなく,それぞれの重心が描く円周である.体積の定理では面積 $A$ と面積重心を,面積の定理では長さ $L$ と線重心を組にして使う.

\begin{bookpoint}
\textbf{POINT}\par
体積は(面積重心が描く円周)$\times$(面積),回転面の面積は(線重心が描く円周)$\times$(長さ)で求める.
\end{bookpoint}

\begin{example}{半円板の回転}
$a>0$,\,$b>0$ とする.上半円板 $x^2+y^2\leq a^2$,\,$y\geq0$ を,直線 $y=-b$ のまわりに1回転してできる立体の体積と全表面積を求めよ.
\end{example}

\paragraphbf{方針}
回転半径は,重心の高さに $b$ を加えたものである.体積には半円板の面積重心を使う.表面積は境界を円弧と直径に分け,それぞれの線重心を使う.

\paragraphbf{解答}
体積 $\pi^2a^2b+\frac{4\pi a^3}{3}$,全表面積 $2\pi^2ab+4\pi a^2+4\pi ab$

\paragraphbf{解法}
図\ref{図:半円の二つの重心}のように,回転軸は図形の外にある.

\textbf{体積}\quad 半円板の面積は $\frac{\pi a^2}{2}$,面積重心の高さは式\eqref{式:半円板の重心}より $\frac{4a}{3\pi}$ である.回転半径は $b+\frac{4a}{3\pi}$ なので,
\begin{align*}
V&=2\pi\left(b+\frac{4a}{3\pi}\right)\cdot\frac{\pi a^2}{2}\\
 &=\pi^2a^2b+\frac{4\pi a^3}{3}.
\end{align*}
断面でも確かめておこう.位置 $x$ の断面は,外半径 $b+\sqrt{a^2-x^2}$,内半径 $b$ の円環である.よって
\begin{align*}
V&=\pi\int_{-a}^a\left\{\left(b+\sqrt{a^2-x^2}\right)^2-b^2\right\}\dd x\\
 &=\pi\int_{-a}^a(a^2-x^2)\dd x+2\pi b\int_{-a}^a\sqrt{a^2-x^2}\dd x
\end{align*}
となる.第1項は $\pi\cdot\frac{4a^3}{3}=\frac{4\pi a^3}{3}$ である.第2項の積分は半円板の面積 $\frac{\pi a^2}{2}$ に等しいので,第2項は $\pi^2a^2b$ である.よって結果は一致する.

\begin{figure}[htbp]
\centering
\begin{tikzpicture}[scale=1.25,>=Stealth]
\fill[blue!10](-2,0)arc[start angle=180,end angle=0,radius=2]--cycle;
\draw[blue!75!black,thick](2,0)arc[start angle=0,end angle=180,radius=2];
\draw[thick](-2,0)--(2,0);
\draw[->](-2.5,0)--(2.6,0)node[right]{$x$};
\draw[->](0,-1.5)--(0,2.45)node[above]{$y$};
\draw[orange!80!black,thick](-2.5,-1)--(2.5,-1)node[right]{$y=-b$};
\fill[blue!75!black](0,0.8488)circle(0.045);
\fill[red!75!black](0,1.2732)circle(0.045);
\draw[blue!75!black](0,0.8488)--(2.5,0.6)node[right]{面積重心 $\frac{4a}{3\pi}$};
\draw[red!75!black](0,1.2732)--(2.5,1.6)node[right]{線重心 $\frac{2a}{\pi}$};
\draw[<->](-2.25,-1)--(-2.25,0)node[midway,left]{$b$};
\node[below]at(-2,0){$-a$};\node[below]at(2,0){$a$};
\end{tikzpicture}
\caption{同じ半円でも,板と円弧では重心が違う.回転半径は軸から測る.}\label{図:半円の二つの重心}
\end{figure}

\textbf{全表面積}\quad 立体の表面は,境界の円弧と直径がそれぞれ回ってできる.円弧の長さは $\pi a$,線重心の高さは式\eqref{式:半円弧の重心}より $\frac{2a}{\pi}$ であるから,円弧が作る面の面積は
\begin{align*}
S_{\mathrm{円弧}}
 &=2\pi\left(b+\frac{2a}{\pi}\right)\cdot\pi a\\
 &=2\pi^2ab+4\pi a^2.
\end{align*}
これだけでは全表面積にならない.直径 $-a\leq x\leq a$,\,$y=0$ も,回転すると半径 $b$,軸方向の長さ $2a$ の円筒面を作る.その面積は
\[
S_{\mathrm{直径}}=2\pi b\cdot2a=4\pi ab
\]
である.2つの面は端の円周で接するだけで,面積のある部分では重ならない.よって全表面積は
\[
S=2\pi^2ab+4\pi a^2+4\pi ab
\]
となる.\,$x=\pm a$ では外半径と内半径がともに $b$ で立体の厚みが $0$ になるので,端を塞ぐ円板を加える必要はない.
%
%
\section*{章末問題}
\begin{enumerate}[label=\textbf{\thechapter.\arabic*},leftmargin=*]
\item\label{問:面積と円環}
$y=x^2$ と $y=3x$ に囲まれた部分の面積を求めよ.また,この部分を $x$ 軸のまわりに回転させた立体の体積を求めよ.
\item\label{問:横帯の面積}
$x=y^2$ と $x=3-2y$ に囲まれた部分の面積を,\,$y$ を積分変数として求めよ.
\item\label{問:外半径の交代}
$0\leq x\leq2$,\,$-1\leq y\leq x$ の領域を $x$ 軸のまわりに回した立体の体積を求めよ.外半径が変わる位置を明示せよ.
\item\label{問:円筒と円板}
$0\leq x\leq1$,\,$x^2\leq y\leq1$ の領域を $y$ 軸のまわりに回した立体の体積を,円筒殻と円板の2通りで求めよ.
\item\label{問:弧長の平方完成}
$y=\frac13(x^2+2)^{\frac32}$,\,$0\leq x\leq1$ の曲線の長さを求めよ.
\item\label{問:道のりと平均速度}
$0\leq t\leq3$ における速度が $v(t)=3(t-1)(t-2)$,初期位置が $s(0)=0$ の運動について,終点の位置,道のり,平均速度,平均の速さを求めよ.
\item\label{問:平均値を取る点}
$f(x)=1+x^2$ の $[0,2]$ における平均値を求め,\,$f(c)$ がその平均値と等しくなる $c\in[0,2]$ を求めよ.
\item\label{問:三角形と両パップス}
頂点 $(0,0),(2,0),(2,1)$ の三角形を考える.面積重心の $y$ 座標を積分で求め,パップス＝ギュルダンの定理により,三角形を $x$ 軸のまわりに回した立体の体積を求めよ.また,境界の各線分の線重心を使って,できる立体の全表面積を求めよ.回転軸にある線分の寄与も説明せよ.
\end{enumerate}

\smallskip\noindent 解答は\pageref{解答:第9章}ページ.

%
